\subsection{Elementary properties of measurable functions}

THEOREM 1. --- Let X be a locally compact space, \(\mu\) a measure on X, \((\mathrm{F}_n)\) a sequence of topological spaces, \(\mathbf{F} = \prod_{n} \mathbf{F}_n\) their product. For each index \(n\) , let \(f_n\) be a measurable mapping of X into \(\mathbf{F}_n\) , and let \(f(x) = (f_n(x)) \in \mathbf{F}\) ; then, for every continuous mapping \(u\) of \(f(\mathbf{X})\) into a topological space G, the function \(x \mapsto u(f(x))\) is measurable.

For every compact subset K of X and every number \(\varepsilon > 0\) , there exists a compact set \(K_{1} \subset K\) such that \(|\mu|(\mathrm{K} - \mathrm{K}_{1}) \leqslant \varepsilon\) and such that the restrictions to \(K_{1}\) of all the functions \(f_{n}\) are continuous (No.~1, Prop. 2); it is clear that \(u \circ f\) is continuous on \(K_{1}\) , whence the theorem.

Remarks. --- 1) The theorem does not extend to an arbitrary product of topological spaces (Exer. 1).

\pandocbounded{\includegraphics[keepaspectratio]{images/part001_9d1170b3.jpg}}

\begin{enumerate}
\def\labelenumi{\arabic{enumi})}
\setcounter{enumi}{1}
\tightlist
\item
  If \(f\) is a continuous mapping of \(X\) into itself, and \(g\) is a measurable mapping of \(X\) into \(F\) , then \(g \circ f\) is not necessarily measurable (Exer. 2).
\end{enumerate}

COROLLARY 1. --- The upper envelope and the lower envelope of a finite number of measurable numerical functions (finite or not) are measurable.

For, \(\sup (u,v)\) and \(\inf (u,v)\) are continuous on \(\overline{\mathbf{R}}\times \overline{\mathbf{R}}\)

COROLLARY 2. --- For a numerical function \(f\) (finite or not) to be measurable, it is necessary and sufficient that \(f^{+}\) and \(f^{-}\) be measurable.

The condition is necessary by Cor. 1; it is sufficient because the image A of X in \(\overline{R} \times \overline{R}\) under the mapping \(x \mapsto (f^{+}(x), f^{-}(x))\) does not contain the points \((+\infty, +\infty)\) and \((-\infty, -\infty)\) , consequently the mapping \((u, v) \mapsto u - v\) is continuous on A.

COROLLARY 3. --- If f and g are two measurable mappings of X into a topological vector space F, then \(f + g\) and \(\alpha f\) are measurable ( \(\alpha\) any scalar).

The set of measurable mappings of X into a topological vector space F is thus a vector space.

COROLLARY 4. --- Let F be a vector space of dimension n over R and let \((\mathbf{e}_{k})_{1\leqslant k\leqslant n}\) be a basis of F. In order that a function \(f=\sum_{k=1}^{n}\mathbf{e}_{k}f_{k}\) be measurable, it is necessary and sufficient that each of the numerical functions \(f_{k}\) be measurable.

COROLLARY 5. --- Let F, G, H be three topological vector spaces, and let \((u, v) \to [u \cdot v]\) be a continuous bilinear mapping of \(F \times G\) into H. If f is a measurable mapping of X into F, and g is a measurable mapping of X into G, then \([f \cdot g]\) is a measurable mapping of X into H.

In particular, if f is a measurable mapping of X into a real (resp. complex) normed space F, and g is a measurable mapping of X into R (resp. C), then gf is measurable. If F is a normed algebra and f,g are two measurable mappings of X into F, then fg is measurable.

COROLLARY 6. --- If f is a measurable mapping of X into a normed space F, then the numerical function \(|f|\) is measurable.

